\documentclass[10pt,a4paper]{amsart}
\usepackage[margin=19mm]{geometry}
\usepackage{amsmath,amssymb,mathtools}
\usepackage[scr=rsfs,cal=euler]{mathalfa}
\usepackage{xfrac}
\usepackage[unicode,hidelinks,bookmarks=false]{hyperref}
\newcommand{\ie}{i.e.\ }
\newcommand{\cf}{cf.\ }
\newcommand{\resp}{respectively\ }
\newcommand{\Subsection}{\subsection}
\providecommand{\nobreakdash}{\nobreak\hbox{-}}
\setlength{\emergencystretch}{2em}
\multlinegap0pt
\usepackage{bm}
\usepackage{bbm}
\usepackage{dsfont}
\usepackage{xspace}
\usepackage{cancel}
\usepackage{mathtools}
\usepackage{amsthm}
\makeatletter
\@ifundefined{theorem}{\newtheorem{theorem}{Theorem}}{}
\makeatother
\title[Gibbs measure for mixed spins and mixed types model]{Gibbs measure for mixed spins and mixed types model}
\author{Muzaffar Rahmatullaev, Akbarkhuja Tukhtabaev}
\date{Source: arXiv:2511.01507v1 (CC BY 4.0)}
\begin{document}
\maketitle

\noindent\emph{Source and licence.} Gibbs measure for mixed spins and mixed types model. Version arXiv:2511.01507v1 (CC BY 4.0), distributed under the Creative Commons Attribution 4.0 International licence, as recorded in the arXiv licence metadata for this exact version and shown on the abstract page https://arxiv.org/abs/2511.01507v1. Full rights notice: licenses/arxiv-2511.01507-CC-BY-4.0.txt.\par\smallskip

\noindent\textit{Editorial scope.} Counted unit: the Theorem whose source label is \texttt{comp.eq} in the article (source0.tex lines 235--243, source label \texttt{comp.eq}), delivered together with the article's own complete original proof of it (source0.tex lines 245--367, one \texttt{proof} environment of 123 lines). No mathematics is written, repaired or re-derived: statement and proof are the article's own text line for line. Required context retained uncounted: Ham (lines 192--194); mu (lines 207--210); comcon (lines 223--226). Every definition, display and label of those lines is retained unchanged. Omitted from this package and not redistributed: 1--191, 195--206, 211--222, 227--234, 244--244, 368--746. Nothing inside a retained excerpt is omitted or altered: every retained body is delivered exactly as the article's lines are written, so no omission receipt of any kind accompanies this package. Citations: the retained excerpts carry no citation command at all, so no bibliography entry of the article is transcribed and no reference is redistributed. Reference adaptation: none was needed and none was made; every reference inside the delivered unit resolves inside it, so no unresolved reference or citation is produced. Printed slips kept literal: no printed word, symbol, hypothesis or number of the counted statement or of its proof was altered. Layout and class: the article's own class is \texttt{sn-jnl} with packages including textgreek, amscd, amssymb, amsmath, amsthm, graphicx, caption, subcaption, xcolor, epstopdf; the wrapper is the lane's pinned \texttt{amsart} editorial wrapper at 10pt on A4, which restates the theorem-like environments the retained text needs and adds the article's own macro definitions that the retained text needs (none), with extra wrapper packages (bm, bbm, dsfont, xspace, cancel, mathtools). The delivered numbering is the unit's own. Assets: no retained span contains a figure environment, a graphics-inclusion command, a PDF-page inclusion, a table or a TikZ picture; no image file is copied into this package, the package declares no asset dependency and no image file is redistributed with it. Licence: version arXiv:2511.01507v1 of this article is distributed under the Creative Commons Attribution 4.0 International licence, as recorded in the arXiv licence metadata for this exact version and shown on the abstract page https://arxiv.org/abs/2511.01507v1. A full rights notice is delivered at licenses/arxiv-2511.01507-CC-BY-4.0.txt. Characters: every retained excerpt is pure ASCII, so the delivered engine, XeLaTeX with its default Latin Modern text font, reports no missing character.
\begin{equation}\label{Ham}
H_n(\sigma_n, s_n)=-\alpha \sum_{\langle x,y\rangle \in L_n}J_{I}(x,y)\sigma(x)\sigma(y)-(1-\alpha)\sum_{\langle x,y\rangle \in L_n}J_{P}(x,y)\delta_{s(x)s(y)},
\end{equation}

\begin{equation}\label{mu}
\mu_{n}(\sigma_n, s_n)=Z_{n}^{-1}\exp\left\{-\beta H_n(\sigma_n,s_n)+\sum _{x\in
W_n}h_{\sigma(x),s(x),x}\right\},
\end{equation}

\begin{equation}\label{comcon}
\sum\limits_{(\sigma_{[n]}, s_{[n]})  \in {\Omega _{{W_n}}}} {\mu_{n}((\sigma_{n-1}, s_{n-1}) \vee (\sigma_{[n]}, s_{[n]})) = \mu _{n -
1}(\sigma_{n-1}, s_{n-1})}.
\end{equation}

\begin{theorem}\label{comp.eq} The measures $\mu_{n}(\sigma_n, s_n), n=1,2,...,$ satisfy the compatibility condition \eqref{comcon} if and only if for any $x\in V\setminus\{x_0\}$ the following equation holds:

If $(i,j)\neq(-1,q)$ then
$$z_{i,j,x}=$$
\begin{equation}\label{funceq}
\prod\limits_{y\in S(x)}\frac{\theta_{I}^{\alpha (1-i)}\left(\sum\limits_{v=1}^{q}\left(\theta_{I}^{2\alpha i}z_{1,v,y}+z_{-1,v,y} \right)+\left(\theta_{P}^{1-\alpha}-1\right)\left(\theta_{I}^{2\alpha i}z_{1,j,y}+z_{-1,j,y}\right)\right) }{\sum\limits_{v=1}^{q-1}\left(z_{1,v,y}+\theta_{I}^{2\alpha }(x,y)z_{-1,v,y} \right)+\theta_{P}^{1-\alpha}\left(z_{1,q,y}+\theta_{I}^{2\alpha }(x,y)\right)}\end{equation}

If $(i,j)=(-1,q)$ then  $z_{-1,q,x}=1$.
\end{theorem}

\begin{proof}
\emph{Necessity.} We prove that \eqref{comcon} $\Rightarrow$ \eqref{funceq}.

We first compute some necessary calculation.

1) We rewrite \eqref{Ham}
\begin{eqnarray*}
H_n(\sigma_n, s_n)&=&-\alpha \sum_{\langle x,y\rangle \in L_n}J_{I}(x,y)\sigma(x)\sigma(y)-(1-\alpha)\sum_{\langle x,y\rangle \in L_n}J_{P}(x,y)\delta_{s(x)s(y)}\\[2mm]
&=&-\alpha \sum_{\langle x,y\rangle \in L_{n-1}}J_{I}(x,y)\sigma(x)\sigma(y)
-\alpha\sum_{x\in W_{n-1}}\sum_{y\in S(x)}J_{I}(x,y)\sigma(x)\sigma(y)\\[2mm]
&-&(1-\alpha) \sum_{\langle x,y\rangle \in L_{n-1}}J_{P}(x,y)\delta_{s(x)s(y)}\\[2mm]
&-&(1-\alpha)\sum_{x\in W_{n-1}}\sum_{y\in S(x)}J_{P}(x,y)\delta_{s(x)s(y)}\\[2mm]
&=&H_{n-1}(\sigma_{n-1}, s_{n-1})-\alpha\sum_{x\in W_{n-1}}\sum_{y\in S(x)}J_{I}(x,y)\sigma_{n-1}(x)\sigma_{[n]}(y)\\[2mm]
&-&(1-\alpha)\sum_{x\in W_{n-1}}\sum_{y\in S(x)}J_{P}(x,y)\delta_{s_{n-1}(x)s_{[n]}(y)},\\[2mm]
\end{eqnarray*}
i.e.

$$H_n(\sigma_n, s_n)=H_{n-1}(\sigma_{n-1}, s_{n-1})-\alpha\sum_{x\in W_{n-1}}\sum_{y\in S(x)}J_{I}(x,y)\sigma_{n-1}(x)\sigma_{[n]}(y)$$
\begin{equation}\label{rH}
-(1-\alpha)\sum_{x\in W_{n-1}}\sum_{y\in S(x)}J_{P}(x,y)\delta_{s_{n-1}(x)s_{[n]}(y)}.
\end{equation}
2) We simplify the following expression
\begin{equation}\label{nthg}
\sum _{x\in W_n}h_{\sigma(x),s(x),x}=\sum_{x\in W_{n-1}}\sum_{y\in S(x)}h_{\sigma(y),s(y),y}=\sum_{x\in W_{n-1}}\sum_{y\in S(x)}h_{\sigma_{[n]}(y),s_{[n]}(y),y}.
\end{equation}
3) Putting \eqref{rH} and \eqref{nthg} into \eqref{mu}, we have
\begin{eqnarray*}
\mu_{n}(\sigma_n, s_n)&=&Z_{n}^{-1}\exp\left\{-\beta H_n(\sigma_n,s_n)+\sum _{x\in
W_n}h_{\sigma(x),s(x),x}\right\}\\[2mm]
\end{eqnarray*}
$$=Z_n^{-1}\exp\{-\beta H_{n-1}(\sigma_{n-1},s_{n-1})+$$ $$+\sum_{x\in W_{n-1}}\sum_{y\in S(x)}\left[\beta \alpha  J_{I}(x,y)\sigma_{n-1}(x)\sigma_{[n]}(y)
+\beta (1-\alpha) J_{P}(x,y)\delta_{s_{n-1}(x)s_{[n]}(y)}\right]+$$ $$+\sum_{x\in W_{n-1}}\sum_{y\in S(x)}\left[ h_{\sigma_{[n]}(y),s_{[n]}(y),y}\right]\}.$$

$$\mu_{n}(\sigma_n, s_n)=Z_n^{-1}\exp\{-\beta H_{n-1}(\sigma_{n-1},s_{n-1})\}\cdot$$

$$\cdot\prod_{x\in W_{n-1}}\prod_{y\in S(x)}\exp\{\left[\beta \alpha  J_{I}(x,y)\sigma_{n-1}(x)\sigma_{[n]}(y)
+\beta (1-\alpha) J_{P}(x,y)\delta_{s_{n-1}(x)s_{[n]}(y)}\right]\}$$
\begin{equation}\label{rM}\cdot\prod_{x\in W_{n-1}}\prod_{y\in S(x)}\exp\{ h_{\sigma_{[n]}(y),s_{[n]}(y),y}\}.
\end{equation}
4) \eqref{mu} follows that

$$\mu_{n}(\sigma_{n-1}, s_{n-1})=Z_{n-1}^{-1}\exp\{-\beta H_{n-1}(\sigma_{n-1},s_{n-1})\}$$
\begin{equation}\label{n1}\cdot\prod_{x\in
W_{n-1}}\exp\{h_{\sigma_{n-1}(x),s_{n-1}(x),x}\}.
\end{equation}
5) \eqref{comcon} yields that

 $$\mu _{n -1}(\sigma_{n-1}, s_{n-1})=\sum\limits_{(\sigma_{[n]}, s_{[n]})\in \Omega _{W_n}} \mu
_{n}((\sigma_{n-1}, s_{n-1}) \vee (\sigma_{[n]}, s_{[n]}))$$$$=\sum\limits_{(\sigma_{[n]}, s_{[n]})  \in {\Omega _{W_n}}} \mu
_{n}(\sigma_n,s_n).$$

Putting \eqref{n1} and \eqref{rM} into the last equation, we get
$$\prod _{x\in
W_{n-1}}\exp\{h_{\sigma_{n-1}(x),s_{n-1}(x),x}\}=$$
$$=\frac{Z_{n-1}}{Z_n}\sum\limits_{(\sigma_{[n]}, s_{[n]})  \in {\Omega _{W_n}}} \prod_{x\in W_{n-1}}\prod_{y\in S(x)}\exp[\beta \alpha  J_{I}(x,y)\sigma_{n-1}(x)\sigma_{[n]}(y)$$
\begin{equation}\label{cc3}
+\beta (1-\alpha) J_{P}(x,y)\delta_{s_{n-1}(x)s_{[n]}(y)}+ h_{\sigma_{[n]}(y),s_{[n]}(y),y}].
\end{equation}

Fix $x\in W_{n-1}$, rewrite the last equality for $(\sigma_{n-1},s_{n-1})=(i,j)\in\ \Phi\times\Psi$ and $(\sigma_{n-1},s_{n-1})=(-1,q)$, then dividing each of them to the last one, we have
$$
\prod\limits_{y\in S(x)}\frac{\sum\limits_{(u,v)\in \Phi\times\Psi}\exp\left[\beta \alpha  J_{I}(x,y)iu
+\beta (1-\alpha) J_{P}(x,y)\delta_{jv}+ h_{u,v,y}\right]}{\sum\limits_{(u,v)\in \Phi\times\Psi}\exp\left[-\beta \alpha  J_{I}(x,y)u
+\beta (1-\alpha) J_{P}(x,y)\delta_{qv}+ h_{u,v,y}\right]}$$
\begin{equation}\label{mainone}=\exp\{h_{i,j,x}-h_{-1,q,x}\},
\end{equation}
where $(i,j)\in\Phi\times\Psi$.

Using notations above $\theta_{I}(x,y)=\theta_{I}=\exp\{\beta J_{I}(x,y)\}$, $\theta_{P}(x,y)=\theta_{P}=\exp\{\beta J_{P}(x,y)\}$, the fraction LHS of \eqref{mainone} equals
$$\frac{\sum\limits_{v\in \Psi}\theta_{P}^{(1-\alpha)\delta_{jv}}\theta_{I}^{-\alpha i}\left(\theta_{I}^{2\alpha i}\exp\{h_{1,v,y}\}+\exp\{h_{-1,v,y}\} \right)}{\sum\limits_{v\in \Psi}\theta_{P}^{(1-\alpha)\delta_{qv} }\theta_{I}^{-\alpha }\left( \exp\{ h_{1,v,y}\}+ \theta_{I}^{2\alpha }\exp\{h_{-1,v,y}\} \right)}
=$$
$$=\frac{\theta_{I}^{-\alpha i}\left(\sum\limits_{v\neq j}\left(\theta_{I}^{2\alpha i}e^{h_{1,v,y}}+e^{h_{-1,v,y}} \right)+\theta_{P}^{1-\alpha} \left(\theta_{I}^{2\alpha i}e^{h_{1,j,y}}+e^{h_{-1,j,y}}\right)\right)}{\theta_{I}^{-\alpha}\left(\sum\limits_{v\neq q}\left(e^{h_{1,v,y}}+\theta_{I}^{2\alpha}e^{h_{-1,v,y}} \right)+\theta_{P}^{1-\alpha}\left(e^{h_{1,q,y}}+\theta_{I}^{2\alpha}e^{h_{-1,q,y}}\right)\right) }$$

Again using notation above $z_{i,j,x}=\exp\{h_{i,j,x}-h_{-1,q,y}\}$, $i=-1,1, j=1,2,...,q$, we have  $z_{-1,q,x}=1$ and
\eqref{mainone} follows that
$$z_{i,j,x}=$$
$$
\prod\limits_{y\in S(x)}\frac{\theta_{I}^{\alpha (1-i)}\left(\sum\limits_{v=1}^{q}\left(\theta_{I}^{2\alpha i}z_{1,v,y}+z_{-1,v,y} \right)+\left(\theta_{P}^{1-\alpha}-1\right)\left(\theta_{I}^{2\alpha i}z_{1,j,y}+z_{-1,j,y}\right)\right) }{\sum\limits_{v=1}^{q-1}\left(z_{1,v,y}+\theta_{I}^{2\alpha }z_{-1,v,y} \right)+\theta_{P}^{1-\alpha}\left(z_{1,q,y}+\theta_{I}^{2\alpha }\right)}.
$$

\emph{Sufficiency.} We prove that \eqref{funceq} $\Rightarrow$ \eqref{comcon}. Suppose that \eqref{funceq} holds. It is equivalent to the representation
$$\prod\limits_{y\in S(x)}\sum\limits_{(u,v)\in \Phi\times\Psi}\exp\left[\beta \alpha  J_{I}(x,y)iu
+\beta (1-\alpha) J_{P}(x,y)\delta_{jv}+ h_{u,v,y}\right]=$$
\begin{equation}\label{cc9}
a(x)\exp\{h_{i,j,x}\}
\end{equation}
for some function $a(x)>0, x\in V$.

We rewrite LHS of \eqref{comcon} as

$$\sum\limits_{(\sigma_{[n]}, s_{[n]})  \in {\Omega _{{W_n}}}} \mu_{n}((\sigma_{n-1}, s_{n-1}) \vee (\sigma_{[n]}, s_{[n]})) =Z_n^{-1}\exp\{-\beta H_{n-1}(\sigma_{n-1},s_{n-1})\}\cdot$$
\begin{equation}\label{cc12}
\cdot\prod_{x\in W_{n-1}}\prod_{y\in S(x)}\sum\limits_{(u,v)\in \Phi\times\Psi}\exp\{\beta \alpha  J_{I}\sigma_{n-1}(x)u
+\beta (1-\alpha) J_{P}\delta_{s_{n-1}(x)v}+ h_{u,v,y}\}.
\end{equation}

Substituting \eqref{cc9} into \eqref{cc12} and denoting $A_n(x)=\prod\limits_{x\in W_{n-1}}a(x)$, we get
$$\sum\limits_{(\sigma_{[n]}, s_{[n]})  \in {\Omega _{{W_n}}}} \mu_{n}((\sigma_{n-1}, s_{n-1}) \vee (\sigma_{[n]}, s_{[n]}))$$
\begin{equation}\label{cc13}
  = \frac{A_{n-1}}{Z_n}\exp\{-\beta H_{n-1}(\sigma_{n-1},s_{n-1})\} \prod\limits_{x\in W_{n-1}}\exp\{h_{\sigma_{n-1},s_{n-1},x}\}
\end{equation}
Since $\mu_{n}$ is a probability, we should have
$$\sum\limits_{(\sigma_{n-1}, s_{n-1})\in \Omega_{V_{n-1}}}\sum\limits_{(\sigma_{[n]}, s_{[n]})  \in {\Omega _{{W_n}}}} \mu_{n}((\sigma_{n-1}, s_{n-1}) \vee (\sigma_{[n]}, s_{[n]}))=1.$$
\eqref{cc13} yields
\begin{equation}\label{cc14}
\sum\limits_{(\sigma_{[n]}, s_{[n]})  \in {\Omega _{{W_n}}}} \mu_{n}((\sigma_{n-1}, s_{n-1}) \vee (\sigma_{[n]}, s_{[n]})) = \frac{A_{n-1}}{Z_n}\mu_{n-1}(\sigma_{n-1}, s_{n-1}) Z_{n-1}.
\end{equation}
or

$$1=\sum\limits_{(\sigma_{n-1}, s_{n-1})\in \Omega_{V_{n-1}}}\sum\limits_{(\sigma_{[n]}, s_{[n]})  \in {\Omega _{{W_n}}}} \mu_{n}((\sigma_{n-1}, s_{n-1}) \vee (\sigma_{[n]}, s_{[n]})),$$
\begin{equation}\label{cc15}
1=\sum\limits_{(\sigma_{n-1}, s_{n-1})\in \Omega_{V_{n-1}}}\frac{A_{n-1}}{Z_n}\mu_{n-1}(\sigma_{n-1}, s_{n-1}) Z_{n-1},
\end{equation}
or
\begin{equation}\label{cc16}
Z_n=A_{n-1}Z_{n-1}.
\end{equation}
Substituting \eqref{cc16} into \eqref{cc14}, we have

\begin{equation}\label{cc17}
\sum\limits_{(\sigma_{[n]}, s_{[n]})  \in {\Omega _{{W_n}}}} \mu_{n}((\sigma_{n-1}, s_{n-1}) \vee (\sigma_{[n]}, s_{[n]})) = \mu_{n-1}(\sigma_{n-1}, s_{n-1}).
\end{equation}
\end{proof}

\end{document}
